% Actual closed P141, source head 5cab33a2e239f3e6d90e63e32c54ba0764b1c594.
% Original pipeline closed 2026-10-28 03:36 JST; no interim outcome selection.
\subsection{表現を固定した最終診断}
P141の全3条件を評価した結果を表\ref{tab:p141}に示す．TRとTDOは成功し，費用はそれぞれ332,588，329,781トークンであった．TDOはTRより2,807トークン（0.84\%）少なく，内訳はactorの差2,753，要約の差54である．モデル要求は31から30へ減ったが，API要求は76から78へ増えた．第二差分は採用され，費用が確定した7回のactor入力で使用された．この1課題では表現とobject指示を合わせた変更に小さな改善が観測された．

\begin{table}[t]
\caption{P141の既知第4課題の結果（各1回）．}
\ecaption{Final diagnostic on known task 4 in P141, one run per policy.}
\label{tab:p141}
\centering
\begin{tabular}{llrr}
\hline
方式 & 成否 & 既知token小計 & 不明\\ \hline
TR & 成功 & 332,588 & 0\\
TDO & 成功 & 329,781 & 0\\
EDO & 失敗 & 317,920 & 1\\ \hline
\end{tabular}
\end{table}

同じ差分表現を用いるEDOは失敗し，使用量不明が1件残った．TDOとEDOの発動判定は19モデル呼出し後の境界で初めて分かれ，その境界の実際の送信役割も異なり，両送信の費用は確定していた．境界までの履歴は，検証済みprotocol IDを正規化すると一致したが，厳密一致は成立しなかった．TRとTDOの表現変更前も，同じく正規化一致のみ成立した．これらは保存された入力・観測の対応を示す検査であり，隠れた状態や意味の同一性を保証しない．

したがって，表現を同じにしても発動方式の変更後に成功が維持されない例を得た．ただし，最終比較は事後選定した1課題の記述的診断であり，TDOの採用を確認した実験ではない．0.84\%を一般的な改善率とせず，P140の非発動3課題を加えて4課題で確認したとも扱わない．EDOを含む確定費用差は算出できない．
